\documentclass[a4paper,11pt]{article}
        \usepackage[top=15mm,bottom=18mm,left=16mm,right=16mm]{geometry}
        \usepackage{fontspec}
        \usepackage{xeCJK}
        \setmainfont{Times New Roman}
        \setCJKmainfont{Yu Gothic}
        \setmonofont{Consolas}
        \setCJKmonofont{Yu Gothic}
        \usepackage{booktabs}
        \usepackage{longtable}
        \usepackage{tabularx}
        \usepackage{array}
        \usepackage{enumitem}
        \usepackage{hyperref}
        \usepackage{listings}
        \usepackage{xcolor}
        \usepackage{amsmath}
        \usepackage{amssymb}
        \hypersetup{
          colorlinks=true,
          linkcolor=blue,
          urlcolor=blue,
          pdftitle={live 自動校正ノード},
          pdfauthor={Codex}
        }
        \lstset{
          basicstyle=\ttfamily\small,
          breaklines=true,
          columns=fullflexible,
          keepspaces=true,
          frame=single,
          framerule=0.2pt,
          rulecolor=\color{black},
          showstringspaces=false
        }
        \setlength{\parskip}{0.42em}
        \setlength{\parindent}{1em}
        \setlength{\emergencystretch}{3em}
        \renewcommand{\arraystretch}{1.15}
        \title{24. live 自動校正ノード}
        \author{solar\_ws0129-main}
        \date{2026-07-29}
        \begin{document}
        \maketitle
        \tableofcontents
        \clearpage

        \section{対象}
        \begin{tabularx}{\linewidth}{>{\raggedright\arraybackslash}p{0.18\linewidth} X}
        \toprule
        項目 & 内容 \\
        \midrule
        ファイル & \path|mpc_solarcar/solar_autocal_node.py| \\
        ソースSHA-256 & \path|a98336e789670e85549bcedbb19c1a90b49d293dba6f193a6d91c053ef9abc51| \\
        種別 & Python \\
        区分 & runtime node \\
        役割 & 観測 solar/pack power と planner 予測との差から solar\_gain、drive\_power\_gain、aux\_power\_w を推定 publish する。 \\
        起動文脈 & 運用中に mpc\_node の係数を微調整する補助ノード。 \\
        \bottomrule
        \end{tabularx}

        \section{このファイルを読む理由}
        観測 solar/pack power と planner 予測との差から solar\_gain、drive\_power\_gain、aux\_power\_w を推定 publish する。

        \begin{itemize}[leftmargin=1.6em]
        \item /calib/* topic を publish する。
\item mpc\_node がそれを購読して内部 gain を更新する。
        \end{itemize}

        \section{呼び出し関係}
        \subsection{どこから呼ばれるか}
        \begin{itemize}[leftmargin=1.6em]
        \item \path|mpc_solarcar/live_launch.py|
        \end{itemize}

        \subsection{次に読むと理解が進むファイル}
        \begin{itemize}[leftmargin=1.6em]
        \item \path|mpc_solarcar/solar_autocal_logic.py|
\item \path|mpc_solarcar/mpc_node.py|
        \end{itemize}

        \subsection{同一リポジトリ内の主要依存}
        \begin{itemize}[leftmargin=1.6em]
        \item \path|mpc_solarcar/solar_autocal_logic.py|
        \end{itemize}

        \section{ソース冒頭の把握}
        モジュール docstring または先頭意図の要約:

        モジュール docstring は明示されていない。

        \subsection{代表コード断片}
        \begin{lstlisting}[language=Python]
return max(float(lo), min(float(hi), float(value)))


class SolarAutocalNode(Node):
    def __init__(self) -> None:
        super().__init__("solar_autocal_node")
        self.declare_parameter("publish_period_sec", 30.0)
        self.declare_parameter("stationary_speed_kmh", 2.0)
        self.declare_parameter("drive_speed_kmh", 25.0)
        self.declare_parameter("day_ghi_threshold", 150.0)
        self.declare_parameter("alpha", 0.2)
        self.declare_parameter("solar_gain_init", 1.0)
        self.declare_parameter("drive_power_gain_init", 1.0)
        self.declare_parameter("aux_power_w_init", 8.0)
        self.declare_parameter("solar_gain_min", 0.5)
        self.declare_parameter("solar_gain_max", 1.5)
        self.declare_parameter("drive_power_gain_min", 0.7)
        self.declare_parameter("drive_power_gain_max", 1.4)
        self.declare_parameter("aux_power_w_min", 0.0)
        self.declare_parameter("aux_power_w_max", 300.0)

        self.period = max(1.0, float(self.get_parameter("publish_period_sec").value))
\end{lstlisting}


        \section{主要構造の読み取り}
        主要クラスは SolarAutocalNode。 主要関数は main。 ROS パラメータ宣言は 14 件。 ROS I/O は publisher 4 件、subscription 6 件。

        \subsection{主要クラス・関数}
            \begin{longtable}{p{0.07\linewidth} p{0.16\linewidth} p{0.25\linewidth} p{0.40\linewidth}}
            \toprule
            行 & 種別 & 名前 & 補足 \\
            \midrule
            17 & class & \path|SolarAutocalNode| & bases: Node \\
13 & function & \path|_clamp| & args: value, lo, hi \\
18 & function & \path|__init__| & args: self \\
74 & function & \path|_on_speed| & args: self, msg \\
77 & function & \path|_on_current| & args: self, msg \\
82 & function & \path|_on_voltage| & args: self, msg \\
87 & function & \path|_on_solar| & args: self, msg \\
90 & function & \path|_on_env| & args: self, msg \\
95 & function & \path|_on_metrics| & args: self, msg \\
102 & function & \path|_ema| & args: self, current, estimate, lo, hi \\
106 & function & \path|_step| & args: self \\
153 & function & \path|main| &  \\
            \bottomrule
            \end{longtable}

        \subsection{ファイルを上から読んだときの定義順}
            \begin{longtable}{p{0.07\linewidth} p{0.84\linewidth}}
            \toprule
            行 & 内容 \\
            \midrule
            13 & 関数 \_clamp を定義する。 \\
17 & クラス SolarAutocalNode を定義する。 \\
153 & 関数 main を定義する。 \\
            \bottomrule
            \end{longtable}

        \subsection{import 群の詳細}
            \begin{longtable}{p{0.07\linewidth} p{0.33\linewidth} p{0.50\linewidth}}
            \toprule
            行 & import 文 & このファイルで読む理由 \\
            \midrule
            1 & \path|from __future__ import annotations| & 型ヒントの遅延評価を有効にし、自己参照型や forward reference を安全に扱うため。 このファイル内での使用位置は少ないか、間接利用である。 \\
3 & \path|import math| & Python標準のスカラー数値関数と定数を使うため。実際に参照する名前と行は使用位置欄で確認する。 このファイル内での主な使用位置は L50, L51, L52, L53, L54, L55, L69, L70, ...。 \\
5 & \path|import rclpy| & ROS 2 Python ノードとして起動・spin するため。 このファイル内での主な使用位置は L154, L156, L158。 \\
6 & \path|from rclpy.node import Node| & ROS 2 ノード本体の基底クラスとして使うため。 このファイル内での主な使用位置は L17。 \\
8 & \path|from std_msgs.msg import Float32, Float32MultiArray, String| & Float32、Bool、Stringなどの標準ROS messageで値をpublishまたはsubscribeするため。 このファイル内での主な使用位置は L57, L58, L59, L60, L62, L63, L64, L65, ...。 \\
10 & \path|from .solar_autocal_logic import daytime_stationary_aux_estimate| & 自動校正ロジック関数群 から daytime\_stationary\_aux\_estimate を読み込み、このファイルの内部処理を分担させるため。 実体ファイルは mpc\_solarcar/solar\_autocal\_logic.py。 このファイル内での主な使用位置は L110。 \\
            \bottomrule
            \end{longtable}

        \section{このファイルを読む前に必要な基礎知識}
        次の章は、構文やROS用語を既知と仮定しないための説明である。

        \subsection{プログラム、プロセス、メモリ、オブジェクトを区別する}
ソースファイルはディスク上の文字列であり、それ自体は走っていない。Python実行可能プログラムがソースを読み、OSがその実行を一つのプロセスとして管理する。プロセスには仮想メモリ、開いているファイル、スレッド、終了コードなどが対応する。


\begin{lstlisting}
ソースファイル -> Pythonインタプリタが読み込む -> OS上のプロセス -> Pythonオブジェクトをメモリに生成 -> 関数やcallbackを実行
\end{lstlisting}


「メモリ上に生成する」とは、実行中プロセスが使う記憶領域に、その値の型、属性、参照関係を表す実体を用意することである。変数は実体そのものというより、そのオブジェクトを指す名前として理解するとPythonの代入が読みやすい。


\begin{lstlisting}[language=python]
node = MPCNode()
alias = node
# nodeとaliasは同じオブジェクトを参照する。
\end{lstlisting}


一つのプロセスには複数スレッドを持たせられる。スレッドは同じプロセスのメモリを共有するため、受信callbackと最適化callbackが同じself.zを書き換える場合は実行順と排他を検討する必要がある。

\paragraph{根拠資料}
\begin{itemize}[leftmargin=1.6em]
\item \href{https://docs.python.org/3/tutorial/classes.html}{Python公式チュートリアル: Classes}
\end{itemize}

\subsection{名前、代入、型、参照}
Pythonの代入文は、右辺を先に評価し、その結果のオブジェクトへ左辺の名前を結び付ける。`x = f()`では、まずfを呼び、その戻り値が得られてからxが更新される。


`float(...)`、`int(...)`、`bool(...)`は型名を呼び出して値を変換する。外部CSV、YAML、ROS parameterから来た値は型が期待どおりとは限らないため、このリポジトリでは明示変換が多い。


`None`は値が無いことを示す単一のオブジェクト、`math.nan`は浮動小数点値ではあるが有効な数値ではないことを示す。両者は用途が異なるため、`is None`と`math.isfinite`を使い分ける。

\paragraph{根拠資料}
\begin{itemize}[leftmargin=1.6em]
\item \href{https://docs.python.org/3/tutorial/classes.html}{Python公式チュートリアル: Classes}
\end{itemize}

\subsection{関数、引数、戻り値、スコープ}
`def`は関数本体をその場で実行する文ではない。関数オブジェクトを作り、その名前を現在の名前空間へ登録する。`f`は関数そのもの、`f()`はその関数を呼んだ結果である。


仮引数は関数定義側の受け取り名、実引数は呼出側から渡す値である。位置引数、キーワード引数、既定値付き引数、`*args`、`**kwargs`は渡し方が異なる。


\begin{lstlisting}[language=python]
def clip_speed(value, minimum=0.0, maximum=110.0):
    return min(maximum, max(minimum, value))

v = clip_speed(120.0, maximum=90.0)
\end{lstlisting}


関数内で作った通常の名前はローカルスコープに属する。メソッドが`self.z`を更新すればオブジェクトの状態が残るが、単に`z`を更新しただけならその呼出しのローカル値である。

\paragraph{根拠資料}
\begin{itemize}[leftmargin=1.6em]
\item \href{https://docs.python.org/3/tutorial/classes.html}{Python公式チュートリアル: Classes}
\end{itemize}

\subsection{module、package、import、console\_scripts}
Pythonファイルはmoduleとして読み込める。複数moduleをディレクトリにまとめたものがpackageである。`from .model import SolarCarModel`の先頭の点は、現在と同じpackage内のmodel moduleを指す相対importである。


import時にはファイルのトップレベル文が上から一度実行される。`def`や`class`は関数・クラスを登録するが、その本体の通常処理は呼び出すまで走らない。


setuptoolsの`console\_scripts`は、端末で使う実行可能名と`package.module:function`を対応付けるインストール時メタデータである。生成される実行用ラッパーはmoduleをimportし、指定関数を呼び、戻り値を終了コードとして扱う小さな入口である。


\begin{lstlisting}
ROS launchのexecutable名 -> install済みconsole script -> mpc_solarcar.mpc_nodeをimport -> main()を呼ぶ
\end{lstlisting}

\paragraph{根拠資料}
\begin{itemize}[leftmargin=1.6em]
\item \href{https://setuptools.pypa.io/en/latest/userguide/entry_point.html}{setuptools公式: Entry Points / Console Scripts}
\item \href{https://docs.python.org/3/tutorial/classes.html}{Python公式チュートリアル: Classes}
\end{itemize}

\subsection{例外、try/finally、with、資源解放}
例外は通常の戻り値とは別経路で異常を呼出元へ伝える。`try/except`は想定した異常を処理し、`finally`は成功・失敗にかかわらず後始末を行う。


`with open(...) as f:`はcontext managerを使い、ブロックを出るとファイルを閉じる。CSVやログの破損を避けるため、開いた資源の所有者と閉じる場所を明確にする。


`except Exception: pass`はノードを止めない利点がある一方、入力異常を隠して原因追跡を難しくする。安全に関係する値では、少なくとも頻度制限付き警告、異常カウンタ、fallback状態のpublishを検討する。



\subsection{class、独自クラス、インスタンス、self、継承}
クラスは、保持するデータと、そのデータを扱う関数を一つの型としてまとめる設計単位である。「独自クラス」とはPythonやROS 2が最初から用意した型ではなく、このプロジェクトが目的に合わせて定義した新しい型を指す。


`class MPCNode(Node)`は、rclpyのNodeを基底クラスとして継承し、Nodeが持つpublisher、subscription、timer、parameterなどの機能にMPC固有機能を追加する。丸括弧内は関数引数ではなく基底クラス指定である。


`MPCNode()`はクラスオブジェクトを呼び出して新しいインスタンスを作る。Pythonは新しい実体を用意し、その実体を第1引数selfとして`\_\_init\_\_`へ渡す。`self`は予約語ではなく慣習的な名前だが、この慣習を守ることでコードの意味が共有される。


\begin{lstlisting}[language=python]
class Counter:
    def __init__(self):
        self.value = 0

    def increment(self):
        self.value += 1

counter = Counter()
counter.increment()
\end{lstlisting}


`super().\_\_init\_\_(...)`は継承元の初期化処理を呼ぶ。MPCNodeではこれを省くとROSノードとして必要な内部実体が作られず、`create\_publisher`などを正しく使えない。

\paragraph{根拠資料}
\begin{itemize}[leftmargin=1.6em]
\item \href{https://docs.python.org/3/tutorial/classes.html}{Python公式チュートリアル: Classes}
\item \href{https://docs.ros.org/en/humble/Tutorials/Beginner-CLI-Tools/Understanding-ROS2-Nodes/Understanding-ROS2-Nodes.html}{ROS 2 Humble公式: Understanding nodes}
\end{itemize}

\subsection{先頭アンダースコア、\_\_init\_\_、名前付けの作法}
`self.\_step\_solar`の先頭1個のアンダースコアは「クラス内部で使う実装詳細」という慣習であり、アクセスを強制禁止する機構ではない。外部コードから呼べるが、公開APIとして依存しない意思表示である。


`\_\_init\_\_`の前後2個のアンダースコアはPythonが定めた特殊メソッド名である。クラス生成時に自動で呼ばれる。任意の名前に前後2個を付けて独自仕様を作ることは避ける。


`\_`で始まるローカル変数は、未使用または外部へ見せない意図を示す場合がある。例えば`\_base\_csv`は戻り値として受け取る必要はあるが、その後の処理では使わないことを読者へ示す。

\paragraph{根拠資料}
\begin{itemize}[leftmargin=1.6em]
\item \href{https://docs.python.org/3/tutorial/classes.html}{Python公式チュートリアル: Classes}
\end{itemize}

\subsection{rclpy、rcl、rmw、DDS/RTPSの層}
rclpyはPython利用者向けROS 2 client libraryである。利用者のNode、publisher、subscriptionなどをPython APIとして提供し、その下で共通C層のrclを利用する。


\begin{lstlisting}
本プロジェクトのPythonコード -> rclpy -> rcl -> rmw -> DDS/RTPS実装 -> ネットワークまたは同一PC内通信
\end{lstlisting}


rclは言語に依存しない共通ROS機能を提供するC API、rmwはROS 2と具体的middleware実装の境界である。DDS/RTPS側が探索、serialize、publish/subscribe、request/replyなどを担う。executorの実行モデルはrclだけで完結せずclient library側にも実装される。

\paragraph{根拠資料}
\begin{itemize}[leftmargin=1.6em]
\item \href{https://docs.ros.org/en/humble/Concepts/About-Internal-Interfaces.html}{ROS 2 Humble公式: Internal ROS 2 interfaces}
\end{itemize}

\subsection{ROS graph、Node、topic、service、action、parameter}
ROS graphは、実行中Nodeと、それらが持つpublisher、subscription、service、actionなどの接続関係である。Pythonクラス、プロセス、ROS graph上のNode名は関連するが同一物ではない。


topicは継続データ向けの非同期一方向publish/subscribe、serviceは短いrequest/response、actionは時間のかかる目標へfeedback、cancel、resultを持たせる。parameterはNodeごとの設定値である。


publisherが送った時点とsubscription callbackが実行される時点は同じとは限らない。通信遅延、QoS queue、executor待ちを挟むため、センサ値にはtimestampとfreshness判定が必要になる。

\paragraph{根拠資料}
\begin{itemize}[leftmargin=1.6em]
\item \href{https://docs.ros.org/en/humble/Tutorials/Beginner-CLI-Tools/Understanding-ROS2-Nodes/Understanding-ROS2-Nodes.html}{ROS 2 Humble公式: Understanding nodes}
\item \href{https://docs.ros.org/en/humble/Concepts/Basic/Interfaces-Topics-Services-Actions.html}{ROS 2 Humble公式: Topics, Services, Actions}
\item \href{https://docs.ros.org/en/humble/Concepts/Basic/About-Parameters.html}{ROS 2 Humble公式: Parameters}
\end{itemize}

\subsection{callback、timer、Executor、spin、Callback Group}
callbackは、メッセージ受信、timer満了、service requestなどのイベントが成立した後でExecutorから呼ばれる関数である。登録時に`self.\_on\_speed`と括弧なしで渡すのは、今実行せず後で呼ぶ関数オブジェクトを渡すためである。


Executorは実行可能になったcallbackを見つけ、Callback Groupの条件を確認して実行する。`spin()`は終了要求まで待機・dispatchを続けるため、通常運転中はmainの次の行へ戻らない。


MutuallyExclusiveCallbackGroupは同じgroup内のcallbackを同時実行させない。別group間はMultiThreadedExecutorで並行実行し得る。groupを分けただけでは共有属性への完全な排他にはならないため、複数groupが同じ状態を読む・書く場合は設計確認が必要である。


\begin{lstlisting}
DDS受信またはtimer満了 -> wait setでready -> Executorが選択 -> Callback Groupが許可 -> worker threadがcallback実行 -> 終了後Executorへ戻る
\end{lstlisting}

\paragraph{根拠資料}
\begin{itemize}[leftmargin=1.6em]
\item \href{https://docs.ros.org/en/rolling/Concepts/Intermediate/About-Executors.html}{ROS 2公式: Executors}
\item \href{https://docs.ros.org/en/humble/Concepts/About-Internal-Interfaces.html}{ROS 2 Humble公式: Internal ROS 2 interfaces}
\end{itemize}

\subsection{rqt\_graph、ros2 CLI、rosbag2をいつ使うか}
rqt\_graphは接続関係を見る道具であり、数値の正しさや更新周期までは保証しない。起動直後、topic名変更後、publisherが複数存在する疑いがある時に使う。


\begin{lstlisting}[language=bash]
ros2 node list
ros2 node info /mpc_node
ros2 topic list -t
ros2 topic info -v /planner/speed_cmd
ros2 topic hz /planner/speed_cmd
ros2 topic echo /planner/status
rqt_graph
\end{lstlisting}


rosbag2はtopicメッセージを時系列のまま記録・再生する。通信不具合、freshness、再計画trigger、実車とSILSの差を再現可能にするため、本番前試験では制御入力だけでなく原因となる全telemetry、status、parameter情報を記録する。


\begin{lstlisting}[language=bash]
ros2 bag record -o outputs/bags/preflight \
  /vehicle/s_km /vehicle/speed_kmh /vehicle/batt_soc \
  /vehicle/batt_temp_c /vehicle/batt_current_a /vehicle/batt_voltage_v \
  /planner/upper_speed_cmd /planner/speed_cmd /planner/status

ros2 bag info outputs/bags/preflight
ros2 bag play outputs/bags/preflight --clock
\end{lstlisting}


bag再生時はQoS互換性、simulation time、外部publisherとの二重入力に注意する。実車Nodeを同時に動かす場合はnamespaceまたはremappingで入力源を明確に分離する。

\paragraph{根拠資料}
\begin{itemize}[leftmargin=1.6em]
\item \href{https://docs.ros.org/en/ros2_packages/humble/api/rqt_graph/index.html}{ROS 2 Humble公式: rqt\_graph}
\item \href{https://docs.ros.org/en/humble/Tutorials/Beginner-CLI-Tools/Recording-And-Playing-Back-Data/Recording-And-Playing-Back-Data.html}{ROS 2 Humble公式: Recording and playing back data}
\item \href{https://docs.ros.org/en/humble/Tutorials/Beginner-CLI-Tools/Understanding-ROS2-Nodes/Understanding-ROS2-Nodes.html}{ROS 2 Humble公式: Understanding nodes}
\end{itemize}

\subsection{制御、状態、入力、モデル予測制御MPC}
制御対象の内部を表す状態をx、操作入力をu、外乱・予報をwとすると、離散モデルは`x[k+1] = f(x[k], u[k], w[k])`と書ける。ソーラーカーではSoC、電池温度、距離、速度などが状態候補、速度目標や駆動トルクが入力候補になる。


\[
\min_{u_0,\ldots,u_{N-1}} \sum_{k=0}^{N-1}\ell(x_k,u_k,w_k)+V_f(x_N)
\]

MPCは現在状態からNステップ先まで予測し、目的関数と制約を満たす入力系列を求める。ただし実際に適用するのは通常先頭入力だけで、次回は新しい実測状態から再び解く。これがreceding horizonである。


予測モデル、目的関数、制約、ホライズン、solver、初期値のどれかが変わると答えも変わる。「MPCを使う」だけでは仕様は決まらず、これらを単位付きで追う必要がある。

\paragraph{根拠資料}
\begin{itemize}[leftmargin=1.6em]
\item \href{https://docs.scipy.org/doc/scipy/reference/generated/scipy.optimize.minimize.html}{SciPy公式: scipy.optimize.minimize}
\end{itemize}



        \section{関数・クラスを上から順に解説}
        \subsection{L13 関数 \_clamp}
            \begin{itemize}[leftmargin=1.6em]
              \item 行範囲: L13--L14
              \item 所属: module直下
              \item 定義: \path|_clamp(value: float, lo: float, hi: float) -> float|
              \item docstring: 明示されていない。
              \item このブロックが直接呼ぶ主な関数・メソッド: float, max, min
              \item 戻り値の要点: max(float(lo), min(float(hi), float(value)))
              \item 主なローカル代入名: 特になし。
              \item 読み取る主なインスタンス属性: 特になし。
              \item 更新する主なインスタンス属性: 特になし。
              \item 明示的に送出する例外: 明示的raiseはない。
              \item 制御構造の規模: 条件分岐 0、ループ 0、try 0
            \end{itemize}
            \paragraph{この定義を読むためのPython構文}
            \begin{itemize}[leftmargin=1.6em]
            \item `->`以後は戻り値の型注釈であり、実行時の値を自動変換する命令ではない。
            \end{itemize}
            \paragraph{上から順の処理}
            \begin{enumerate}[leftmargin=1.8em]
            \item max(float(lo), min(float(hi), float(value))) を返す。
            \end{enumerate}
            \paragraph{代表コード断片}
            \begin{lstlisting}[language=Python]
def _clamp(value: float, lo: float, hi: float) -> float:
    return max(float(lo), min(float(hi), float(value)))
\end{lstlisting}

\subsection{L17 クラス SolarAutocalNode}
            \begin{itemize}[leftmargin=1.6em]
              \item 行範囲: L17--L150
              \item 所属: module直下
              \item 定義: \path|SolarAutocalNode(bases=Node)|
              \item docstring: 明示されていない。
              \item このブロックが直接呼ぶ主な関数・メソッド: 特に明示されていない。
              \item 戻り値の要点: 戻り値が無いか、return 文が明示されていない。
              \item 主なローカル代入名: 特になし。
              \item 読み取る主なインスタンス属性: 特になし。
              \item 更新する主なインスタンス属性: 特になし。
              \item 明示的に送出する例外: 明示的raiseはない。
              \item 制御構造の規模: 条件分岐 0、ループ 0、try 0
            \end{itemize}
            \paragraph{この定義を読むためのPython構文}
            \begin{itemize}[leftmargin=1.6em]
            \item class文は新しい型を定義する。丸括弧内は継承する基底クラスである。
            \end{itemize}
            \paragraph{上から順の処理}
            \begin{enumerate}[leftmargin=1.8em]
            \item 関数 \_\_init\_\_ を定義する。
\item 関数 \_on\_speed を定義する。
\item 関数 \_on\_current を定義する。
\item 関数 \_on\_voltage を定義する。
\item 関数 \_on\_solar を定義する。
\item 関数 \_on\_env を定義する。
\item 関数 \_on\_metrics を定義する。
\item 関数 \_ema を定義する。
\item 関数 \_step を定義する。
            \end{enumerate}
            \paragraph{代表コード断片}
            \begin{lstlisting}[language=Python]
class SolarAutocalNode(Node):
    def __init__(self) -> None:
        super().__init__("solar_autocal_node")
        self.declare_parameter("publish_period_sec", 30.0)
        self.declare_parameter("stationary_speed_kmh", 2.0)
        self.declare_parameter("drive_speed_kmh", 25.0)
        self.declare_parameter("day_ghi_threshold", 150.0)
        self.declare_parameter("alpha", 0.2)
        self.declare_parameter("solar_gain_init", 1.0)
        self.declare_parameter("drive_power_gain_init", 1.0)
        self.declare_parameter("aux_power_w_init", 8.0)
        self.declare_parameter("solar_gain_min", 0.5)
        self.declare_parameter("solar_gain_max", 1.5)
        self.declare_parameter("drive_power_gain_min", 0.7)
        self.declare_parameter("drive_power_gain_max", 1.4)
        self.declare_parameter("aux_power_w_min", 0.0)
        self.declare_parameter("aux_power_w_max", 300.0)

        self.period = max(1.0, float(self.get_parameter("publish_period_sec").value))
        self.stationary_speed = float(self.get_parameter("stationary_speed_kmh").value)
        self.drive_speed = float(self.get_parameter("drive_speed_kmh").value)
        self.day_ghi = float(self.get_parameter("day_ghi_threshold").value)
        self.alpha = _clamp(float(self.get_parameter("alpha").value), 0.01, 1.0)
        self.solar_gain = float(self.get_parameter("solar_gain_init").value)
        self.drive_gain = float(self.get_parameter("drive_power_gain_init").value)
        self.aux_power = float(self.get_parameter("aux_power_w_init").value)
        self.solar_gain_min = float(self.get_parameter("solar_gain_min").value)
        self.solar_gain_max = float(self.get_parameter("solar_gain_max").value)
        self.drive_gain_min = float(self.get_parameter("drive_power_gain_min").value)
        self.drive_gain_max = float(self.get_parameter("drive_power_gain_max").value)
        self.aux_min = float(self.get_parameter("aux_power_w_min").value)
        self.aux_max = float(self.get_parameter("aux_power_w_max").value)

        self.speed_kmh = math.nan
        self.pack_power_w = math.nan
...
\end{lstlisting}

\subsection{L18 関数 SolarAutocalNode.\_\_init\_\_}
            \begin{itemize}[leftmargin=1.6em]
              \item 行範囲: L18--L72
              \item 所属: \path|SolarAutocalNode|
              \item 定義: \path|__init__(self) -> None|
              \item docstring: 明示されていない。
              \item このブロックが直接呼ぶ主な関数・メソッド: \_\_init\_\_, \_clamp, create\_publisher, create\_subscription, create\_timer, declare\_parameter, float, get\_logger, get\_parameter, info, max, super
              \item 戻り値の要点: 戻り値が無いか、return 文が明示されていない。
              \item 主なローカル代入名: 特になし。
              \item 読み取る主なインスタンス属性: self.\_on\_current, self.\_on\_env, self.\_on\_metrics, self.\_on\_solar, self.\_on\_speed, self.\_on\_voltage, self.\_step, self.create\_publisher, self.create\_subscription, self.create\_timer, self.declare\_parameter, self.get\_logger, self.get\_parameter, self.period
              \item 更新する主なインスタンス属性: self.\_latest\_current\_a, self.\_latest\_voltage\_v, self.alpha, self.aux\_max, self.aux\_min, self.aux\_power, self.day\_ghi, self.drive\_gain, self.drive\_gain\_max, self.drive\_gain\_min, self.drive\_speed, self.ghi, self.pack\_power\_w, self.period, self.pred\_pack\_w, self.pred\_solar\_w, self.pub\_aux, self.pub\_drive, self.pub\_solar, self.pub\_status, self.solar\_gain, self.solar\_gain\_max, self.solar\_gain\_min, self.solar\_obs\_w
              \item 明示的に送出する例外: 明示的raiseはない。
              \item 制御構造の規模: 条件分岐 0、ループ 0、try 0
            \end{itemize}
            \paragraph{この定義を読むためのPython構文}
            \begin{itemize}[leftmargin=1.6em]
            \item 第1引数selfは、このメソッドを呼んだインスタンス自身を受け取る慣習名である。
\item `->`以後は戻り値の型注釈であり、実行時の値を自動変換する命令ではない。
            \end{itemize}
            \paragraph{上から順の処理}
            \begin{enumerate}[leftmargin=1.8em]
            \item super().\_\_init\_\_(...) を実行する。
\item self.declare\_parameter(...) を実行する。
\item self.declare\_parameter(...) を実行する。
\item self.declare\_parameter(...) を実行する。
\item self.declare\_parameter(...) を実行する。
\item self.declare\_parameter(...) を実行する。
\item self.declare\_parameter(...) を実行する。
\item self.declare\_parameter(...) を実行する。
\item self.declare\_parameter(...) を実行する。
\item self.declare\_parameter(...) を実行する。
\item self.declare\_parameter(...) を実行する。
\item self.declare\_parameter(...) を実行する。
\item self.declare\_parameter(...) を実行する。
\item self.declare\_parameter(...) を実行する。
\item self.declare\_parameter(...) を実行する。
\item self.period に max(1.0, float(self.get\_parameter('publish\_period\_sec').value)) の結果を代入する。
\item self.stationary\_speed に float(self.get\_parameter('stationary\_speed\_kmh').value) の結果を代入する。
\item self.drive\_speed に float(self.get\_parameter('drive\_speed\_kmh').value) の結果を代入する。
\item self.day\_ghi に float(self.get\_parameter('day\_ghi\_threshold').value) の結果を代入する。
\item self.alpha に \_clamp(float(self.get\_parameter('alpha').value), 0.01, 1.0) の結果を代入する。
\item self.solar\_gain に float(self.get\_parameter('solar\_gain\_init').value) の結果を代入する。
\item self.drive\_gain に float(self.get\_parameter('drive\_power\_gain\_init').value) の結果を代入する。
\item self.aux\_power に float(self.get\_parameter('aux\_power\_w\_init').value) の結果を代入する。
\item self.solar\_gain\_min に float(self.get\_parameter('solar\_gain\_min').value) の結果を代入する。
\item self.solar\_gain\_max に float(self.get\_parameter('solar\_gain\_max').value) の結果を代入する。
\item self.drive\_gain\_min に float(self.get\_parameter('drive\_power\_gain\_min').value) の結果を代入する。
\item self.drive\_gain\_max に float(self.get\_parameter('drive\_power\_gain\_max').value) の結果を代入する。
\item self.aux\_min に float(self.get\_parameter('aux\_power\_w\_min').value) の結果を代入する。
\item self.aux\_max に float(self.get\_parameter('aux\_power\_w\_max').value) の結果を代入する。
\item self.speed\_kmh に math.nan の結果を代入する。
\item self.pack\_power\_w に math.nan の結果を代入する。
\item self.solar\_obs\_w に math.nan の結果を代入する。
\item self.ghi に math.nan の結果を代入する。
\item self.pred\_solar\_w に math.nan の結果を代入する。
\item self.pred\_pack\_w に math.nan の結果を代入する。
\item self.pub\_solar に self.create\_publisher(Float32, '/calib/solar\_gain', 10) の結果を代入する。
\item self.pub\_drive に self.create\_publisher(Float32, '/calib/drive\_power\_gain', 10) の結果を代入する。
\item self.pub\_aux に self.create\_publisher(Float32, '/calib/aux\_power\_w', 10) の結果を代入する。
\item self.pub\_status に self.create\_publisher(String, '/calib/status', 10) の結果を代入する。
\item self.create\_subscription(...) を実行する。
\item self.create\_subscription(...) を実行する。
\item self.create\_subscription(...) を実行する。
\item self.create\_subscription(...) を実行する。
\item self.create\_subscription(...) を実行する。
\item self.create\_subscription(...) を実行する。
\item self.\_latest\_current\_a に math.nan の結果を代入する。
\item self.\_latest\_voltage\_v に math.nan の結果を代入する。
\item self.timer に self.create\_timer(self.period, self.\_step) の結果を代入する。
\item self.get\_logger().info(...) を実行する。
            \end{enumerate}
            \paragraph{代表コード断片}
            \begin{lstlisting}[language=Python]
    def __init__(self) -> None:
        super().__init__("solar_autocal_node")
        self.declare_parameter("publish_period_sec", 30.0)
        self.declare_parameter("stationary_speed_kmh", 2.0)
        self.declare_parameter("drive_speed_kmh", 25.0)
        self.declare_parameter("day_ghi_threshold", 150.0)
        self.declare_parameter("alpha", 0.2)
        self.declare_parameter("solar_gain_init", 1.0)
        self.declare_parameter("drive_power_gain_init", 1.0)
        self.declare_parameter("aux_power_w_init", 8.0)
        self.declare_parameter("solar_gain_min", 0.5)
        self.declare_parameter("solar_gain_max", 1.5)
        self.declare_parameter("drive_power_gain_min", 0.7)
        self.declare_parameter("drive_power_gain_max", 1.4)
        self.declare_parameter("aux_power_w_min", 0.0)
        self.declare_parameter("aux_power_w_max", 300.0)

        self.period = max(1.0, float(self.get_parameter("publish_period_sec").value))
        self.stationary_speed = float(self.get_parameter("stationary_speed_kmh").value)
        self.drive_speed = float(self.get_parameter("drive_speed_kmh").value)
        self.day_ghi = float(self.get_parameter("day_ghi_threshold").value)
        self.alpha = _clamp(float(self.get_parameter("alpha").value), 0.01, 1.0)
        self.solar_gain = float(self.get_parameter("solar_gain_init").value)
        self.drive_gain = float(self.get_parameter("drive_power_gain_init").value)
        self.aux_power = float(self.get_parameter("aux_power_w_init").value)
        self.solar_gain_min = float(self.get_parameter("solar_gain_min").value)
        self.solar_gain_max = float(self.get_parameter("solar_gain_max").value)
        self.drive_gain_min = float(self.get_parameter("drive_power_gain_min").value)
        self.drive_gain_max = float(self.get_parameter("drive_power_gain_max").value)
        self.aux_min = float(self.get_parameter("aux_power_w_min").value)
        self.aux_max = float(self.get_parameter("aux_power_w_max").value)

        self.speed_kmh = math.nan
        self.pack_power_w = math.nan
        self.solar_obs_w = math.nan
...
\end{lstlisting}

\subsection{L74 関数 SolarAutocalNode.\_on\_speed}
            \begin{itemize}[leftmargin=1.6em]
              \item 行範囲: L74--L75
              \item 所属: \path|SolarAutocalNode|
              \item 定義: \path|_on_speed(self, msg: Float32) -> None|
              \item docstring: 明示されていない。
              \item このブロックが直接呼ぶ主な関数・メソッド: float
              \item 戻り値の要点: 戻り値が無いか、return 文が明示されていない。
              \item 主なローカル代入名: 特になし。
              \item 読み取る主なインスタンス属性: 特になし。
              \item 更新する主なインスタンス属性: self.speed\_kmh
              \item 明示的に送出する例外: 明示的raiseはない。
              \item 制御構造の規模: 条件分岐 0、ループ 0、try 0
            \end{itemize}
            \paragraph{この定義を読むためのPython構文}
            \begin{itemize}[leftmargin=1.6em]
            \item 第1引数selfは、このメソッドを呼んだインスタンス自身を受け取る慣習名である。
\item `->`以後は戻り値の型注釈であり、実行時の値を自動変換する命令ではない。
            \end{itemize}
            \paragraph{上から順の処理}
            \begin{enumerate}[leftmargin=1.8em]
            \item self.speed\_kmh に float(msg.data) の結果を代入する。
            \end{enumerate}
            \paragraph{代表コード断片}
            \begin{lstlisting}[language=Python]
    def _on_speed(self, msg: Float32) -> None:
        self.speed_kmh = float(msg.data)
\end{lstlisting}

\subsection{L77 関数 SolarAutocalNode.\_on\_current}
            \begin{itemize}[leftmargin=1.6em]
              \item 行範囲: L77--L80
              \item 所属: \path|SolarAutocalNode|
              \item 定義: \path|_on_current(self, msg: Float32) -> None|
              \item docstring: 明示されていない。
              \item このブロックが直接呼ぶ主な関数・メソッド: float, isfinite
              \item 戻り値の要点: 戻り値が無いか、return 文が明示されていない。
              \item 主なローカル代入名: 特になし。
              \item 読み取る主なインスタンス属性: self.\_latest\_current\_a, self.\_latest\_voltage\_v
              \item 更新する主なインスタンス属性: self.\_latest\_current\_a, self.pack\_power\_w
              \item 明示的に送出する例外: 明示的raiseはない。
              \item 制御構造の規模: 条件分岐 1、ループ 0、try 0
            \end{itemize}
            \paragraph{この定義を読むためのPython構文}
            \begin{itemize}[leftmargin=1.6em]
            \item 第1引数selfは、このメソッドを呼んだインスタンス自身を受け取る慣習名である。
\item `->`以後は戻り値の型注釈であり、実行時の値を自動変換する命令ではない。
            \end{itemize}
            \paragraph{上から順の処理}
            \begin{enumerate}[leftmargin=1.8em]
            \item self.\_latest\_current\_a に float(msg.data) の結果を代入する。
\item 条件 math.isfinite(self.\_latest\_voltage\_v) を判定し、真なら内部処理を行う。
\item   self.pack\_power\_w に self.\_latest\_current\_a * self.\_latest\_voltage\_v の結果を代入する。
            \end{enumerate}
            \paragraph{代表コード断片}
            \begin{lstlisting}[language=Python]
    def _on_current(self, msg: Float32) -> None:
        self._latest_current_a = float(msg.data)
        if math.isfinite(self._latest_voltage_v):
            self.pack_power_w = self._latest_current_a * self._latest_voltage_v
\end{lstlisting}

\subsection{L82 関数 SolarAutocalNode.\_on\_voltage}
            \begin{itemize}[leftmargin=1.6em]
              \item 行範囲: L82--L85
              \item 所属: \path|SolarAutocalNode|
              \item 定義: \path|_on_voltage(self, msg: Float32) -> None|
              \item docstring: 明示されていない。
              \item このブロックが直接呼ぶ主な関数・メソッド: float, isfinite
              \item 戻り値の要点: 戻り値が無いか、return 文が明示されていない。
              \item 主なローカル代入名: 特になし。
              \item 読み取る主なインスタンス属性: self.\_latest\_current\_a, self.\_latest\_voltage\_v
              \item 更新する主なインスタンス属性: self.\_latest\_voltage\_v, self.pack\_power\_w
              \item 明示的に送出する例外: 明示的raiseはない。
              \item 制御構造の規模: 条件分岐 1、ループ 0、try 0
            \end{itemize}
            \paragraph{この定義を読むためのPython構文}
            \begin{itemize}[leftmargin=1.6em]
            \item 第1引数selfは、このメソッドを呼んだインスタンス自身を受け取る慣習名である。
\item `->`以後は戻り値の型注釈であり、実行時の値を自動変換する命令ではない。
            \end{itemize}
            \paragraph{上から順の処理}
            \begin{enumerate}[leftmargin=1.8em]
            \item self.\_latest\_voltage\_v に float(msg.data) の結果を代入する。
\item 条件 math.isfinite(self.\_latest\_current\_a) を判定し、真なら内部処理を行う。
\item   self.pack\_power\_w に self.\_latest\_current\_a * self.\_latest\_voltage\_v の結果を代入する。
            \end{enumerate}
            \paragraph{代表コード断片}
            \begin{lstlisting}[language=Python]
    def _on_voltage(self, msg: Float32) -> None:
        self._latest_voltage_v = float(msg.data)
        if math.isfinite(self._latest_current_a):
            self.pack_power_w = self._latest_current_a * self._latest_voltage_v
\end{lstlisting}

\subsection{L87 関数 SolarAutocalNode.\_on\_solar}
            \begin{itemize}[leftmargin=1.6em]
              \item 行範囲: L87--L88
              \item 所属: \path|SolarAutocalNode|
              \item 定義: \path|_on_solar(self, msg: Float32) -> None|
              \item docstring: 明示されていない。
              \item このブロックが直接呼ぶ主な関数・メソッド: float
              \item 戻り値の要点: 戻り値が無いか、return 文が明示されていない。
              \item 主なローカル代入名: 特になし。
              \item 読み取る主なインスタンス属性: 特になし。
              \item 更新する主なインスタンス属性: self.solar\_obs\_w
              \item 明示的に送出する例外: 明示的raiseはない。
              \item 制御構造の規模: 条件分岐 0、ループ 0、try 0
            \end{itemize}
            \paragraph{この定義を読むためのPython構文}
            \begin{itemize}[leftmargin=1.6em]
            \item 第1引数selfは、このメソッドを呼んだインスタンス自身を受け取る慣習名である。
\item `->`以後は戻り値の型注釈であり、実行時の値を自動変換する命令ではない。
            \end{itemize}
            \paragraph{上から順の処理}
            \begin{enumerate}[leftmargin=1.8em]
            \item self.solar\_obs\_w に float(msg.data) の結果を代入する。
            \end{enumerate}
            \paragraph{代表コード断片}
            \begin{lstlisting}[language=Python]
    def _on_solar(self, msg: Float32) -> None:
        self.solar_obs_w = float(msg.data)
\end{lstlisting}

\subsection{L90 関数 SolarAutocalNode.\_on\_env}
            \begin{itemize}[leftmargin=1.6em]
              \item 行範囲: L90--L93
              \item 所属: \path|SolarAutocalNode|
              \item 定義: \path|_on_env(self, msg: Float32MultiArray) -> None|
              \item docstring: 明示されていない。
              \item このブロックが直接呼ぶ主な関数・メソッド: float, len, list
              \item 戻り値の要点: 戻り値が無いか、return 文が明示されていない。
              \item 主なローカル代入名: data
              \item 読み取る主なインスタンス属性: 特になし。
              \item 更新する主なインスタンス属性: self.ghi
              \item 明示的に送出する例外: 明示的raiseはない。
              \item 制御構造の規模: 条件分岐 1、ループ 0、try 0
            \end{itemize}
            \paragraph{この定義を読むためのPython構文}
            \begin{itemize}[leftmargin=1.6em]
            \item 第1引数selfは、このメソッドを呼んだインスタンス自身を受け取る慣習名である。
\item `->`以後は戻り値の型注釈であり、実行時の値を自動変換する命令ではない。
            \end{itemize}
            \paragraph{上から順の処理}
            \begin{enumerate}[leftmargin=1.8em]
            \item data に list(msg.data) の結果を代入する。
\item 条件 len(data) >= 1 を判定し、真なら内部処理を行う。
\item   self.ghi に float(data[0]) の結果を代入する。
            \end{enumerate}
            \paragraph{代表コード断片}
            \begin{lstlisting}[language=Python]
    def _on_env(self, msg: Float32MultiArray) -> None:
        data = list(msg.data)
        if len(data) >= 1:
            self.ghi = float(data[0])
\end{lstlisting}

\subsection{L95 関数 SolarAutocalNode.\_on\_metrics}
            \begin{itemize}[leftmargin=1.6em]
              \item 行範囲: L95--L100
              \item 所属: \path|SolarAutocalNode|
              \item 定義: \path|_on_metrics(self, msg: Float32MultiArray) -> None|
              \item docstring: 明示されていない。
              \item このブロックが直接呼ぶ主な関数・メソッド: float, len, list
              \item 戻り値の要点: 戻り値が無いか、return 文が明示されていない。
              \item 主なローカル代入名: data
              \item 読み取る主なインスタンス属性: 特になし。
              \item 更新する主なインスタンス属性: self.pred\_pack\_w, self.pred\_solar\_w
              \item 明示的に送出する例外: 明示的raiseはない。
              \item 制御構造の規模: 条件分岐 2、ループ 0、try 0
            \end{itemize}
            \paragraph{この定義を読むためのPython構文}
            \begin{itemize}[leftmargin=1.6em]
            \item 第1引数selfは、このメソッドを呼んだインスタンス自身を受け取る慣習名である。
\item `->`以後は戻り値の型注釈であり、実行時の値を自動変換する命令ではない。
            \end{itemize}
            \paragraph{上から順の処理}
            \begin{enumerate}[leftmargin=1.8em]
            \item data に list(msg.data) の結果を代入する。
\item 条件 len(data) >= 6 を判定し、真なら内部処理を行う。
\item   self.pred\_solar\_w に float(data[5]) の結果を代入する。
\item 条件 len(data) >= 9 を判定し、真なら内部処理を行う。
\item   self.pred\_pack\_w に float(data[8]) の結果を代入する。
            \end{enumerate}
            \paragraph{代表コード断片}
            \begin{lstlisting}[language=Python]
    def _on_metrics(self, msg: Float32MultiArray) -> None:
        data = list(msg.data)
        if len(data) >= 6:
            self.pred_solar_w = float(data[5])
        if len(data) >= 9:
            self.pred_pack_w = float(data[8])
\end{lstlisting}

\subsection{L102 関数 SolarAutocalNode.\_ema}
            \begin{itemize}[leftmargin=1.6em]
              \item 行範囲: L102--L104
              \item 所属: \path|SolarAutocalNode|
              \item 定義: \path|_ema(self, current: float, estimate: float, lo: float, hi: float) -> float|
              \item docstring: 明示されていない。
              \item このブロックが直接呼ぶ主な関数・メソッド: \_clamp, float
              \item 戻り値の要点: \_clamp(blended, lo, hi)
              \item 主なローカル代入名: blended
              \item 読み取る主なインスタンス属性: self.alpha
              \item 更新する主なインスタンス属性: 特になし。
              \item 明示的に送出する例外: 明示的raiseはない。
              \item 制御構造の規模: 条件分岐 0、ループ 0、try 0
            \end{itemize}
            \paragraph{この定義を読むためのPython構文}
            \begin{itemize}[leftmargin=1.6em]
            \item 第1引数selfは、このメソッドを呼んだインスタンス自身を受け取る慣習名である。
\item `->`以後は戻り値の型注釈であり、実行時の値を自動変換する命令ではない。
            \end{itemize}
            \paragraph{上から順の処理}
            \begin{enumerate}[leftmargin=1.8em]
            \item blended に (1.0 - self.alpha) * float(current) + self.alpha * float(estimate) の結果を代入する。
\item \_clamp(blended, lo, hi) を返す。
            \end{enumerate}
            \paragraph{代表コード断片}
            \begin{lstlisting}[language=Python]
    def _ema(self, current: float, estimate: float, lo: float, hi: float) -> float:
        blended = (1.0 - self.alpha) * float(current) + self.alpha * float(estimate)
        return _clamp(blended, lo, hi)
\end{lstlisting}

\subsection{L106 関数 SolarAutocalNode.\_step}
            \begin{itemize}[leftmargin=1.6em]
              \item 行範囲: L106--L150
              \item 所属: \path|SolarAutocalNode|
              \item 定義: \path|_step(self) -> None|
              \item docstring: 明示されていない。
              \item このブロックが直接呼ぶ主な関数・メソッド: Float32, String, \_ema, abs, daytime\_stationary\_aux\_estimate, float, isfinite, max, publish
              \item 戻り値の要点: 戻り値が無いか、return 文が明示されていない。
              \item 主なローカル代入名: aux\_estimate, gain\_est, status
              \item 読み取る主なインスタンス属性: self.\_ema, self.aux\_max, self.aux\_min, self.aux\_power, self.day\_ghi, self.drive\_gain, self.drive\_gain\_max, self.drive\_gain\_min, self.drive\_speed, self.ghi, self.pack\_power\_w, self.pred\_pack\_w, self.pred\_solar\_w, self.pub\_aux, self.pub\_drive, self.pub\_solar, self.pub\_status, self.solar\_gain, self.solar\_gain\_max, self.solar\_gain\_min, self.solar\_obs\_w, self.speed\_kmh, self.stationary\_speed
              \item 更新する主なインスタンス属性: self.aux\_power, self.drive\_gain, self.solar\_gain
              \item 明示的に送出する例外: 明示的raiseはない。
              \item 制御構造の規模: 条件分岐 3、ループ 0、try 0
            \end{itemize}
            \paragraph{この定義を読むためのPython構文}
            \begin{itemize}[leftmargin=1.6em]
            \item 第1引数selfは、このメソッドを呼んだインスタンス自身を受け取る慣習名である。
\item `->`以後は戻り値の型注釈であり、実行時の値を自動変換する命令ではない。
            \end{itemize}
            \paragraph{上から順の処理}
            \begin{enumerate}[leftmargin=1.8em]
            \item aux\_estimate に daytime\_stationary\_aux\_estimate(ghi\_wm2=self.ghi, day\_ghi\_threshold\_wm2=self.day\_ghi, speed\_kmh=self.speed\_kmh, stationary\_speed\_kmh=self.stationary\_speed, pack\_power\_w=self.pack\_power\_w, solar\_power\_w=self.solar\_obs\_w) の結果を代入する。
\item 条件 aux\_estimate is not None を判定し、真なら内部処理を行う。
\item   self.aux\_power に self.\_ema(self.aux\_power, aux\_estimate, self.aux\_min, self.aux\_max) の結果を代入する。
\item 条件 math.isfinite(self.ghi) and self.ghi >= self.day\_ghi and math.isfinite(self.solar\_obs\_w) and math.isfinite(self.pred\_solar\_w) and (abs(self.pred\_solar\_w) >= 20.0) を判定し、真なら内部処理を行う。
\item   gain\_est に self.solar\_obs\_w / self.pred\_solar\_w の結果を代入する。
\item   self.solar\_gain に self.\_ema(self.solar\_gain, gain\_est, self.solar\_gain\_min, self.solar\_gain\_max) の結果を代入する。
\item 条件 math.isfinite(self.speed\_kmh) and self.speed\_kmh >= self.drive\_speed and math.isfinite(self.pack\_power\_w) and math.isfinite(self.pred\_pack\_w) and (abs(self.pred\_pack\_w) >= 50.0) and (self.pack\_power\_w > 0.0) を判定し、真なら内部処理を行う。
\item   gain\_est に self.pack\_power\_w / max(1.0, self.pred\_pack\_w) の結果を代入する。
\item   self.drive\_gain に self.\_ema(self.drive\_gain, gain\_est, self.drive\_gain\_min, self.drive\_gain\_max) の結果を代入する。
\item self.pub\_solar.publish(...) を実行する。
\item self.pub\_drive.publish(...) を実行する。
\item self.pub\_aux.publish(...) を実行する。
\item status に f'solar\_gain=\{self.solar\_gain:.3f\} drive\_power\_gain=\{self.drive\_gain:.3f\} aux\_power\_w=\{self.aux\_power:.1f\}' の結果を代入する。
\item self.pub\_status.publish(...) を実行する。
            \end{enumerate}
            \paragraph{代表コード断片}
            \begin{lstlisting}[language=Python]
    def _step(self) -> None:
        # The team switches auxiliaries off at night, so a night stationary
        # sample is not an estimate of the daytime auxiliary load.  At a
        # daytime stop P_pack = P_aux - P_pv, hence P_aux = P_pack + P_pv.
        aux_estimate = daytime_stationary_aux_estimate(
            ghi_wm2=self.ghi,
            day_ghi_threshold_wm2=self.day_ghi,
            speed_kmh=self.speed_kmh,
            stationary_speed_kmh=self.stationary_speed,
            pack_power_w=self.pack_power_w,
            solar_power_w=self.solar_obs_w,
        )
        if aux_estimate is not None:
            self.aux_power = self._ema(self.aux_power, aux_estimate, self.aux_min, self.aux_max)

        if (
            math.isfinite(self.ghi)
            and self.ghi >= self.day_ghi
            and math.isfinite(self.solar_obs_w)
            and math.isfinite(self.pred_solar_w)
            and abs(self.pred_solar_w) >= 20.0
        ):
            gain_est = self.solar_obs_w / self.pred_solar_w
            self.solar_gain = self._ema(self.solar_gain, gain_est, self.solar_gain_min, self.solar_gain_max)

        if (
            math.isfinite(self.speed_kmh)
            and self.speed_kmh >= self.drive_speed
            and math.isfinite(self.pack_power_w)
            and math.isfinite(self.pred_pack_w)
            and abs(self.pred_pack_w) >= 50.0
            and self.pack_power_w > 0.0
        ):
            gain_est = self.pack_power_w / max(1.0, self.pred_pack_w)
            self.drive_gain = self._ema(self.drive_gain, gain_est, self.drive_gain_min, self.drive_gain_max)
...
\end{lstlisting}

\subsection{L153 関数 main}
            \begin{itemize}[leftmargin=1.6em]
              \item 行範囲: L153--L158
              \item 所属: module直下
              \item 定義: \path|main() -> None|
              \item docstring: 明示されていない。
              \item このブロックが直接呼ぶ主な関数・メソッド: SolarAutocalNode, destroy\_node, init, shutdown, spin
              \item 戻り値の要点: 戻り値が無いか、return 文が明示されていない。
              \item 主なローカル代入名: node
              \item 読み取る主なインスタンス属性: 特になし。
              \item 更新する主なインスタンス属性: 特になし。
              \item 明示的に送出する例外: 明示的raiseはない。
              \item 制御構造の規模: 条件分岐 0、ループ 0、try 0
            \end{itemize}
            \paragraph{この定義を読むためのPython構文}
            \begin{itemize}[leftmargin=1.6em]
            \item `->`以後は戻り値の型注釈であり、実行時の値を自動変換する命令ではない。
            \end{itemize}
            \paragraph{上から順の処理}
            \begin{enumerate}[leftmargin=1.8em]
            \item rclpy.init(...) を実行する。
\item node に SolarAutocalNode() の結果を代入する。
\item rclpy.spin(...) を実行する。
\item node.destroy\_node(...) を実行する。
\item rclpy.shutdown(...) を実行する。
            \end{enumerate}
            \paragraph{代表コード断片}
            \begin{lstlisting}[language=Python]
def main() -> None:
    rclpy.init()
    node = SolarAutocalNode()
    rclpy.spin(node)
    node.destroy_node()
    rclpy.shutdown()
\end{lstlisting}

        \subsection{設定値と入力パラメータ}
            \begin{longtable}{p{0.07\linewidth} p{0.35\linewidth} p{0.45\linewidth}}
            \toprule
            行 & 名前 & 既定値・補足 \\
            \midrule
            20 & \path|publish_period_sec| & 30.0 \\
21 & \path|stationary_speed_kmh| & 2.0 \\
22 & \path|drive_speed_kmh| & 25.0 \\
23 & \path|day_ghi_threshold| & 150.0 \\
24 & \path|alpha| & 0.2 \\
25 & \path|solar_gain_init| & 1.0 \\
26 & \path|drive_power_gain_init| & 1.0 \\
27 & \path|aux_power_w_init| & 8.0 \\
28 & \path|solar_gain_min| & 0.5 \\
29 & \path|solar_gain_max| & 1.5 \\
30 & \path|drive_power_gain_min| & 0.7 \\
31 & \path|drive_power_gain_max| & 1.4 \\
32 & \path|aux_power_w_min| & 0.0 \\
33 & \path|aux_power_w_max| & 300.0 \\
            \bottomrule
            \end{longtable}

        \subsection{ROS topic I/O}
            \begin{longtable}{p{0.07\linewidth} p{0.18\linewidth} p{0.34\linewidth} p{0.28\linewidth}}
            \toprule
            行 & 種別 & topic & callback / 補足 \\
            \midrule
            57 & publisher & \path|/calib/solar_gain| & publisher \\
58 & publisher & \path|/calib/drive_power_gain| & publisher \\
59 & publisher & \path|/calib/aux_power_w| & publisher \\
60 & publisher & \path|/calib/status| & publisher \\
62 & subscription & \path|/vehicle/speed_kmh| & self.\_on\_speed \\
63 & subscription & \path|/vehicle/batt_current_a| & self.\_on\_current \\
64 & subscription & \path|/vehicle/batt_voltage_v| & self.\_on\_voltage \\
65 & subscription & \path|/vehicle/solar_power_w| & self.\_on\_solar \\
66 & subscription & \path|/planner/env| & self.\_on\_env \\
67 & subscription & \path|/planner/metrics| & self.\_on\_metrics \\
            \bottomrule
            \end{longtable}







        \section{処理の流れ}
        \begin{enumerate}[leftmargin=1.7em]
        \item 初期化時に設定値や入力パスを読み込む。
\item publisher / subscription / timer を準備する。
\item timer callback 周期で主処理を進める。
        \end{enumerate}

        \section{このファイルの位置づけ}
        この資料群では、\path|mpc_solarcar/solar_autocal_node.py| を
        \textbf{runtime node}の中核として扱う。
        したがって、ここで扱う処理は単独で完結せず、
        前段の profile / launch / telemetry と後段の planner / logger / report へ繋がる。

        \end{document}
